\documentclass[10pt,a4paper,twocolumn]{article}
\usepackage{subfiles}

% ==============================================================================
% banq-preamble.tex
% Shared LaTeX preamble for all XPower Banq papers and the unified bundle.
% Loaded via % ==============================================================================
% banq-preamble.tex
% Shared LaTeX preamble for all XPower Banq papers and the unified bundle.
% Loaded via % ==============================================================================
% banq-preamble.tex
% Shared LaTeX preamble for all XPower Banq papers and the unified bundle.
% Loaded via \input{../shared/banq-preamble} from each child paper and from
% the master bundle root. Identical to the original per-paper preambles.
% ==============================================================================

\usepackage[utf8]{inputenc}
\usepackage[T1]{fontenc}
\usepackage{newunicodechar}
\usepackage{libertinus}
\usepackage[margin=2cm]{geometry}
\usepackage{amsmath,amssymb,amsthm}
\usepackage{graphicx}
\usepackage{xcolor}
\usepackage[hyphens]{url}
\usepackage{hyperref}
\usepackage{booktabs}
\usepackage{tabularx}
\usepackage{multirow}
\usepackage{enumitem}
\usepackage{tikz}
\usepackage{pgfplots}
\usepackage{listings}
\usepackage{caption}
\usepackage{subcaption}
\usepackage{fancyhdr}
\usepackage{float}
\usepackage{placeins}  % For \FloatBarrier
\usepackage{tocloft}   % Adjustable TOC number-box widths
\usepackage{multicol}  % Multi-column layout (used for the bundle's master TOC)
\usepackage{twemojis}  % Twemoji glyphs (pdflatex-compatible color emoji)

% Float parameters to prevent overlapping
\setcounter{topnumber}{2}
\setcounter{bottomnumber}{2}
\setcounter{totalnumber}{4}
\renewcommand{\topfraction}{0.85}
\renewcommand{\bottomfraction}{0.85}
\renewcommand{\textfraction}{0.15}
\renewcommand{\floatpagefraction}{0.7}

% Prevent line breaks inside inline math (fixes broken superscripts/sqrt)
\relpenalty=10000
\binoppenalty=10000

\pgfplotsset{compat=1.18}
\usetikzlibrary{shapes,arrows,positioning,calc,patterns}

% ==============================================================================
% Colors and Styling
% ==============================================================================
\definecolor{codeblue}{RGB}{0,102,204}
\definecolor{codepurple}{RGB}{128,0,128}
\definecolor{codegreen}{RGB}{0,128,0}
\definecolor{codegray}{RGB}{128,128,128}
\definecolor{primaryblue}{RGB}{41,98,255}
\definecolor{warningred}{RGB}{220,53,69}
\definecolor{successgreen}{RGB}{40,167,69}

\hypersetup{
    colorlinks=true,
    linkcolor=primaryblue,
    citecolor=primaryblue,
    urlcolor=primaryblue
}

\lstset{
    basicstyle=\ttfamily\scriptsize,
    keywordstyle=\color{codeblue}\bfseries,
    commentstyle=\color{codegreen},
    stringstyle=\color{codepurple},
    numbers=left,
    numberstyle=\tiny\color{codegray},
    numbersep=4pt,
    xleftmargin=2em,
    framexleftmargin=2em,
    breaklines=true,
    frame=single,
    captionpos=b
}
\lstdefinelanguage{Solidity}{
    keywords={function, returns, return, if, else, for, while, do, break, continue, throw, require, revert, assert, public, private, internal, external, view, pure, payable, memory, storage, calldata, constant, immutable, virtual, override, modifier, event, emit, struct, enum, mapping, contract, interface, library, abstract, is, new, delete, true, false, uint256, uint128, uint64, uint32, uint16, uint8, uint240, uint224, int256, int128, int64, int32, int8, address, bool, bytes, bytes32, string, unchecked},
    sensitive=true,
    morecomment=[l]{//},
    morecomment=[s]{/*}{*/},
    morestring=[b]",
    morestring=[b]'
}

% ==============================================================================
% Theorem Environments
% ==============================================================================
\theoremstyle{definition}
\newtheorem{definition}{Definition}[section]
\newtheorem{property}{Property}[section]
\theoremstyle{plain}
\newtheorem{theorem}{Theorem}[section]
\newtheorem{lemma}[theorem]{Lemma}
\newtheorem{corollary}[theorem]{Corollary}
\theoremstyle{remark}
\newtheorem{remark}{Remark}[section]
 from each child paper and from
% the master bundle root. Identical to the original per-paper preambles.
% ==============================================================================

\usepackage[utf8]{inputenc}
\usepackage[T1]{fontenc}
\usepackage{newunicodechar}
\usepackage{libertinus}
\usepackage[margin=2cm]{geometry}
\usepackage{amsmath,amssymb,amsthm}
\usepackage{graphicx}
\usepackage{xcolor}
\usepackage[hyphens]{url}
\usepackage{hyperref}
\usepackage{booktabs}
\usepackage{tabularx}
\usepackage{multirow}
\usepackage{enumitem}
\usepackage{tikz}
\usepackage{pgfplots}
\usepackage{listings}
\usepackage{caption}
\usepackage{subcaption}
\usepackage{fancyhdr}
\usepackage{float}
\usepackage{placeins}  % For \FloatBarrier
\usepackage{tocloft}   % Adjustable TOC number-box widths
\usepackage{multicol}  % Multi-column layout (used for the bundle's master TOC)
\usepackage{twemojis}  % Twemoji glyphs (pdflatex-compatible color emoji)

% Float parameters to prevent overlapping
\setcounter{topnumber}{2}
\setcounter{bottomnumber}{2}
\setcounter{totalnumber}{4}
\renewcommand{\topfraction}{0.85}
\renewcommand{\bottomfraction}{0.85}
\renewcommand{\textfraction}{0.15}
\renewcommand{\floatpagefraction}{0.7}

% Prevent line breaks inside inline math (fixes broken superscripts/sqrt)
\relpenalty=10000
\binoppenalty=10000

\pgfplotsset{compat=1.18}
\usetikzlibrary{shapes,arrows,positioning,calc,patterns}

% ==============================================================================
% Colors and Styling
% ==============================================================================
\definecolor{codeblue}{RGB}{0,102,204}
\definecolor{codepurple}{RGB}{128,0,128}
\definecolor{codegreen}{RGB}{0,128,0}
\definecolor{codegray}{RGB}{128,128,128}
\definecolor{primaryblue}{RGB}{41,98,255}
\definecolor{warningred}{RGB}{220,53,69}
\definecolor{successgreen}{RGB}{40,167,69}

\hypersetup{
    colorlinks=true,
    linkcolor=primaryblue,
    citecolor=primaryblue,
    urlcolor=primaryblue
}

\lstset{
    basicstyle=\ttfamily\scriptsize,
    keywordstyle=\color{codeblue}\bfseries,
    commentstyle=\color{codegreen},
    stringstyle=\color{codepurple},
    numbers=left,
    numberstyle=\tiny\color{codegray},
    numbersep=4pt,
    xleftmargin=2em,
    framexleftmargin=2em,
    breaklines=true,
    frame=single,
    captionpos=b
}
\lstdefinelanguage{Solidity}{
    keywords={function, returns, return, if, else, for, while, do, break, continue, throw, require, revert, assert, public, private, internal, external, view, pure, payable, memory, storage, calldata, constant, immutable, virtual, override, modifier, event, emit, struct, enum, mapping, contract, interface, library, abstract, is, new, delete, true, false, uint256, uint128, uint64, uint32, uint16, uint8, uint240, uint224, int256, int128, int64, int32, int8, address, bool, bytes, bytes32, string, unchecked},
    sensitive=true,
    morecomment=[l]{//},
    morecomment=[s]{/*}{*/},
    morestring=[b]",
    morestring=[b]'
}

% ==============================================================================
% Theorem Environments
% ==============================================================================
\theoremstyle{definition}
\newtheorem{definition}{Definition}[section]
\newtheorem{property}{Property}[section]
\theoremstyle{plain}
\newtheorem{theorem}{Theorem}[section]
\newtheorem{lemma}[theorem]{Lemma}
\newtheorem{corollary}[theorem]{Corollary}
\theoremstyle{remark}
\newtheorem{remark}{Remark}[section]
 from each child paper and from
% the master bundle root. Identical to the original per-paper preambles.
% ==============================================================================

\usepackage[utf8]{inputenc}
\usepackage[T1]{fontenc}
\usepackage{newunicodechar}
\usepackage{libertinus}
\usepackage[margin=2cm]{geometry}
\usepackage{amsmath,amssymb,amsthm}
\usepackage{graphicx}
\usepackage{xcolor}
\usepackage[hyphens]{url}
\usepackage{hyperref}
\usepackage{booktabs}
\usepackage{tabularx}
\usepackage{multirow}
\usepackage{enumitem}
\usepackage{tikz}
\usepackage{pgfplots}
\usepackage{listings}
\usepackage{caption}
\usepackage{subcaption}
\usepackage{fancyhdr}
\usepackage{float}
\usepackage{placeins}  % For \FloatBarrier
\usepackage{tocloft}   % Adjustable TOC number-box widths
\usepackage{multicol}  % Multi-column layout (used for the bundle's master TOC)
\usepackage{twemojis}  % Twemoji glyphs (pdflatex-compatible color emoji)

% Float parameters to prevent overlapping
\setcounter{topnumber}{2}
\setcounter{bottomnumber}{2}
\setcounter{totalnumber}{4}
\renewcommand{\topfraction}{0.85}
\renewcommand{\bottomfraction}{0.85}
\renewcommand{\textfraction}{0.15}
\renewcommand{\floatpagefraction}{0.7}

% Prevent line breaks inside inline math (fixes broken superscripts/sqrt)
\relpenalty=10000
\binoppenalty=10000

\pgfplotsset{compat=1.18}
\usetikzlibrary{shapes,arrows,positioning,calc,patterns}

% ==============================================================================
% Colors and Styling
% ==============================================================================
\definecolor{codeblue}{RGB}{0,102,204}
\definecolor{codepurple}{RGB}{128,0,128}
\definecolor{codegreen}{RGB}{0,128,0}
\definecolor{codegray}{RGB}{128,128,128}
\definecolor{primaryblue}{RGB}{41,98,255}
\definecolor{warningred}{RGB}{220,53,69}
\definecolor{successgreen}{RGB}{40,167,69}

\hypersetup{
    colorlinks=true,
    linkcolor=primaryblue,
    citecolor=primaryblue,
    urlcolor=primaryblue
}

\lstset{
    basicstyle=\ttfamily\scriptsize,
    keywordstyle=\color{codeblue}\bfseries,
    commentstyle=\color{codegreen},
    stringstyle=\color{codepurple},
    numbers=left,
    numberstyle=\tiny\color{codegray},
    numbersep=4pt,
    xleftmargin=2em,
    framexleftmargin=2em,
    breaklines=true,
    frame=single,
    captionpos=b
}
\lstdefinelanguage{Solidity}{
    keywords={function, returns, return, if, else, for, while, do, break, continue, throw, require, revert, assert, public, private, internal, external, view, pure, payable, memory, storage, calldata, constant, immutable, virtual, override, modifier, event, emit, struct, enum, mapping, contract, interface, library, abstract, is, new, delete, true, false, uint256, uint128, uint64, uint32, uint16, uint8, uint240, uint224, int256, int128, int64, int32, int8, address, bool, bytes, bytes32, string, unchecked},
    sensitive=true,
    morecomment=[l]{//},
    morecomment=[s]{/*}{*/},
    morestring=[b]",
    morestring=[b]'
}

% ==============================================================================
% Theorem Environments
% ==============================================================================
\theoremstyle{definition}
\newtheorem{definition}{Definition}[section]
\newtheorem{property}{Property}[section]
\theoremstyle{plain}
\newtheorem{theorem}{Theorem}[section]
\newtheorem{lemma}[theorem]{Lemma}
\newtheorem{corollary}[theorem]{Corollary}
\theoremstyle{remark}
\newtheorem{remark}{Remark}[section]

% ==============================================================================
% banq-meta.tex
% Shared author/version metadata for all XPower Banq papers.
% ==============================================================================

\newcommand{\banqAuthor}{\textsc{Karun The Ritch}}
\newcommand{\banqAuthorEmail}{ktr@xpowerbanq.com}
\newcommand{\banqAuthorURL}{https://www.xpowerbanq.com}
\newcommand{\banqAuthorContact}{%
    \href{mailto:\banqAuthorEmail}{\texttt{\banqAuthorEmail}}\quad\url{\banqAuthorURL}%
}
\newcommand{\banqArxivId}{preprint:2605~~[q-fin.CP]}
\newcommand{\banqArxivDate}{1 May 2026}

% ==============================================================================
% banq-bundle.tex
% Defines the \ifbundle conditional that distinguishes a standalone child build
% from the unified bundle build, and the \xref macro that produces a real
% internal hyperlink in bundle mode and a citation-with-text in standalone mode.
%
% Standalone child papers leave \bundlefalse (the default).
% The master bundle root sets \bundletrue before \begin{document}.
% ==============================================================================

\newif\ifbundle
\bundlefalse

% \xref{<bibkey>}{<label>}{<text>}
% - In bundle mode: produces a clickable in-document hyperlink to the label.
% - In standalone mode: prints the text followed by a citation to the bibkey.
% Note: defined after \bundletrue/\bundlefalse so it picks up the right branch.
% Children must call this from inside \AtBeginDocument to ensure \ifbundle is
% finalized; alternatively the macro can test \ifbundle at call site.
\newcommand{\xref}[3]{%
  \ifbundle
    \hyperref[#2]{#3}%
  \else
    #3~\cite{#1}%
  \fi
}

% \xrefshort{<bibkey>}{<label>}{<text>}
% Same as \xref but suppresses the citation entirely in standalone mode (used
% when the surrounding prose already cites the source).
\newcommand{\xrefshort}[3]{%
  \ifbundle
    \hyperref[#2]{#3}%
  \else
    #3%
  \fi
}

% \banqEnableBundleNumbering
% Resets the section counter at every \part and prefixes \thesection with the
% part's Roman numeral, so headings read "I.1 Introduction", "II.1 Introduction",
% etc. Called explicitly from the master root after \bundletrue (see
% 000.banq-all/banq.tex). Not called by standalone children.
%
% Parts that contain MULTIPLE sub-papers must additionally call \banqSubpart
% (defined below) before each child \subfile, which appends a per-paper letter
% suffix (a, b, ...) to the part numeral and resets the section counter so each
% sub-paper restarts from 1. Without this, sub-papers within the same part
% share a single section sequence, which causes the TOC to conflate identically
% titled sections (e.g. two "II.1 Introduction" entries for banq-lck/banq-log).
%
% The outer \makeatletter/\makeatother brackets here are critical: the macro
% body contains \@addtoreset (an @ control sequence) which must be parseable
% as a single token at \newcommand-definition time, not at call time.
\newcommand{\banqSubpartLetter}{}
% Boolean tracking whether we are currently inside a sub-paper of a shared part
% (i.e. \banqSubpart has been called since the last \part). Used by \banqPaper
% to decide whether to demote section bookmark levels. \ifx-comparing
% \banqSubpartLetter against \@empty doesn't work because \newcommand produces
% \long macros while \@empty is short — \ifx always returns false.
\newif\ifbanqSharedPart
\makeatletter
\newcommand{\banqEnableBundleNumbering}{%
  \@addtoreset{section}{part}%
  \renewcommand{\thesection}{\Roman{part}\banqSubpartLetter.\,\arabic{section}}%
  % Make hyperref anchor names unique across sub-papers that share a \part.
  % Without this, banq-lck and banq-log both emit section.2.1, section.2.2, …
  % which corrupts the PDF outline tree (external viewers truncate at Part II).
  \renewcommand{\theHsection}{\Roman{part}\banqSubpartLetter.\arabic{section}}%
  \renewcommand{\theHsubsection}{\theHsection.\arabic{subsection}}%
  \renewcommand{\theHsubsubsection}{\theHsubsection.\arabic{subsubsection}}%
  \renewcommand{\theHparagraph}{\theHsubsubsection.\arabic{paragraph}}%
  % Hook \part so the subpart-letter and shared-part flag both reset when a new
  % part begins; otherwise state set in part N would leak into part N+1.
  \let\banqOriginalPart\part
  \renewcommand{\part}{%
    \renewcommand{\banqSubpartLetter}{}%
    \banqSharedPartfalse
    \banqOriginalPart
  }%
  % The "I.\,10" prefix is wider than the default \numberline box. Widen the
  % TOC number-box widths so two-digit section numbers don't overlap titles.
  \setlength{\cftsecnumwidth}{3.5em}%
  \setlength{\cftsubsecnumwidth}{4.2em}%
  \setlength{\cftsubsubsecnumwidth}{5.2em}%
}
\makeatother

% \banqSubpart{<letter>}
% Marks the start of a sub-paper inside a part that contains multiple papers.
% Sets the per-paper letter suffix appended to the part numeral in \thesection
% (so headings read "IIa.1 Introduction", "IIb.1 Introduction", etc.) and
% resets the section counter so the sub-paper's own numbering restarts at 1.
% Pass an empty argument (\banqSubpart{}) to clear the suffix mid-part if
% needed; the suffix also auto-clears at every \part boundary (see hook above).
\newcommand{\banqSubpart}[1]{%
  \renewcommand{\banqSubpartLetter}{#1}%
  \setcounter{section}{0}%
  \banqSharedParttrue
}

% \banq@RestoreLevels / \banq@DemoteLevels
% Adjust hyperref's per-sectioning bookmark levels. When two sub-papers share a
% \part (banq-lck "IIa" and banq-log "IIb" inside Part~II), each paper's
% sections should nest under its own divider bookmark in the PDF outline,
% rather than appearing as siblings of it under the part. Demoting
% \toclevel@section from 1 to 2 (and the deeper levels accordingly) achieves
% that nesting; \banqPaper calls Restore before emitting its own divider so the
% divider always lands at level 1, then calls Demote afterwards (only when in a
% shared part) so subsequent sections sit at level 2 = child of the divider.
% Affects only the PDF outline depths; the printed TOC layout is unchanged.
\makeatletter
\newcommand{\banq@RestoreLevels}{%
  \renewcommand{\toclevel@section}{1}%
  \renewcommand{\toclevel@subsection}{2}%
  \renewcommand{\toclevel@subsubsection}{3}%
  \renewcommand{\toclevel@paragraph}{4}%
}
\newcommand{\banq@DemoteLevels}{%
  \renewcommand{\toclevel@section}{2}%
  \renewcommand{\toclevel@subsection}{3}%
  \renewcommand{\toclevel@subsubsection}{4}%
  \renewcommand{\toclevel@paragraph}{5}%
}
\makeatother

% Storage for the current part's twemoji name. \banqPartWithEmoji sets it
% (overwritten at each part boundary). \banqPaper reads it without clearing,
% so multiple sub-papers within the same part (e.g. banq-lck and banq-log
% in Part~II) share the same emoji on their respective cover pages.
\makeatletter
\def\banq@partemoji{}
\newif\ifbanq@hasemoji
\banq@hasemojifalse
% \banqSetSubpartEmoji{<twemoji-name>}
% Override the part's default emoji for the next sub-paper cover. Used inside
% parts that contain multiple papers when each sub-paper deserves its own
% cover icon (e.g. banq-lck takes "lock" while banq-log keeps the part-level
% "gear"). Encapsulates the @-internal \gdef so the master document body need
% not enter \makeatletter scope.
\newcommand{\banqSetSubpartEmoji}[1]{\gdef\banq@partemoji{#1}}
\makeatother

% \banqForceRecto
% Page-break helper that flushes outstanding output and ensures the next
% content-bearing page lands on an odd (recto) page in duplex bindings.
% The article class is loaded without `twoside` so \cleardoublepage degenerates
% to \clearpage and will not pad; we check the page counter explicitly and
% insert a blank filler when the upcoming page is even. \value{page} reflects
% the next page number after \clearpage, before any content has been emitted
% on it. Call sites: front-matter (TOC start, body start) and \banqPaper.
\newcommand{\banqForceRecto}{%
  \clearpage
  \ifodd\value{page}\else\hbox{}\thispagestyle{empty}\newpage\fi
}

% \banqEnsureEvenPageCount
% Final-pass parity helper: guarantees the document ends with an even total
% page count, so duplex-bound copies have a back cover on the verso side and
% any subsequent insert (e.g. a colophon, a bound-in errata sheet) starts on
% a recto. After \clearpage, \value{page} holds the *next* page number, so
% the current page count equals \value{page}-1; the count is even iff
% \value{page} is odd, which is exactly the invariant \banqForceRecto already
% enforces. The two operations are semantically distinct (recto-start vs.
% even-total) but compute identical pad logic, so we alias them. Call once,
% immediately before \end{document}, in the bundle root.
\let\banqEnsureEvenPageCount\banqForceRecto

% \banqPartWithEmoji{<title>}{<twemoji-name>}
% Bumps the part counter, stashes the title for the part-level TOC entry,
% and stashes the emoji for the cover page(s) that follow (see \banqPaper).
% No physical divider page is emitted — the cover that immediately follows
% is the visible boundary between parts. The bundle's section-numbering
% scheme (I.1, IIa.1, IIb.1, III.1, ...) continues to work because
% \refstepcounter drives the same counter \part would have, and
% \@addtoreset (set up by \banqEnableBundleNumbering) still resets section
% at the part boundary.
%
% The TOC entry itself is written by \banqPaper at the start of the cover
% page rather than here, so the deferred page-number capture lands on the
% cover (the visible part boundary) rather than on whatever page the build
% happens to be on between subfiles.
\makeatletter
\def\banq@parttitle{}
\newif\ifbanq@parttocpending
\banq@parttocpendingfalse
\newcommand{\banqPartWithEmoji}[2]{%
  \refstepcounter{part}%
  \renewcommand{\banqSubpartLetter}{}%
  \banqSharedPartfalse
  \gdef\banq@partemoji{#2}%
  \global\banq@hasemojitrue
  \gdef\banq@parttitle{#1}%
  \global\banq@parttocpendingtrue
}
\makeatother

% \banqAbstract{<abstract text>}{<keywords>}
% Single source of truth for each child's abstract + keywords list, with
% mode-aware rendering:
%   - Standalone: emits a normal \begin{abstract}...\end{abstract} block
%     followed by a "Keywords:" line (the previous \ifbundle\else ... \fi
%     boilerplate each child used to inline).
%   - Bundle: stashes the content into module-level token lists and sets
%     \ifbanq@hasabstract; \banqPaper then renders it on the verso page that
%     follows the cover (the page that was previously left blank).
% Defined as \long so abstracts containing paragraph breaks are accepted,
% though in practice all six abstracts are single-paragraph.
\makeatletter
\def\banq@abstracttext{}
\def\banq@keywordstext{}
\newif\ifbanq@hasabstract
\banq@hasabstractfalse
\long\def\banqAbstract#1#2{%
  \ifbundle
    \long\gdef\banq@abstracttext{#1}%
    \long\gdef\banq@keywordstext{#2}%
    \global\banq@hasabstracttrue
  \else
    \begin{abstract}#1\end{abstract}\par\medskip\noindent
    \textbf{Keywords:} #2\par
  \fi
}
\makeatother

% \banqPaper{<title>}{<subtitle>}
% Paper-level divider page printed at the top of each child's body when in
% bundle mode. Sits below the \part page in the visual hierarchy and provides
% a clear transition between consecutive papers within the same part (e.g.
% banq-lck -> banq-log inside Part~II). A no-op in standalone builds.
%
% Wrapped in \makeatletter so the macro body can reference \banq@RestoreLevels,
% \banq@DemoteLevels, and \@empty as single tokens rather than \banq + @… .
\makeatletter
\newcommand{\banqPaper}[2]{%
  \ifbundle
    % Force the cover page to land on an odd (recto) page; see \banqForceRecto.
    \banqForceRecto
    \onecolumn
    \thispagestyle{empty}
    % If a \banqPartWithEmoji call is pending, write its TOC + PDF-outline
    % entry now so the captured page number is the cover page (the visible
    % part boundary), not whatever blank pad the build happened to land on.
    \ifbanq@parttocpending
      \phantomsection
      \addcontentsline{toc}{part}{\thepart\hspace{1em}\banq@parttitle}%
      \global\banq@parttocpendingfalse
    \fi
    \vspace*{6em}
    \begin{center}
      {\Large\color{codegray} \textsc{XPower Banq: Part \Roman{part}\banqSubpartLetter}}\\[0.6em]
      \rule{0.4\textwidth}{0.4pt}\\[1.2em]
      {\fontsize{22}{26}\selectfont\bfseries #1}\\[0.8em]
      {\large\color{codegray} #2}
    \end{center}
    \vspace*{\fill}
    \ifbanq@hasemoji
      \begin{center}
        \twemoji[height=8cm]{\banq@partemoji}
      \end{center}
    \fi
    \vspace*{\fill}
    \clearpage
    % Verso page after the cover. If the child supplied an abstract via
    % \banqAbstract, render it here (still in onecolumn, unstyled body width
    % so the prose breathes). Otherwise leave the page blank — used for the
    % "References & Glossary" cover, which has no per-paper abstract since
    % the bundle's front-matter abstract already summarises the whole volume.
    \thispagestyle{empty}
    \ifbanq@hasabstract
      \vspace*{\fill}
      \begin{center}{\large\bfseries Abstract}\end{center}
      \medskip
      \begingroup
        \small
        \noindent\banq@abstracttext\par
        \medskip
        \noindent\textbf{Keywords:} \banq@keywordstext\par
      \endgroup
      \vspace*{\fill}
      \global\banq@hasabstractfalse
    \else
      \hbox{}%
    \fi
    \newpage
    \twocolumn
    % Register the divider in the TOC and PDF outline as a child of the part.
    % Restore section bookmark levels first so this divider always lands at
    % level 1 (regardless of any prior demotion left over from a sibling
    % sub-paper), then conditionally demote so subsequent sections nest under
    % this divider when we're inside a shared part.
    \banq@RestoreLevels
    \phantomsection
    \addcontentsline{toc}{section}{\protect\numberline{\Roman{part}\banqSubpartLetter.}#1}%
    \ifbanqSharedPart
      \banq@DemoteLevels
    \fi
  \fi
}
\makeatother


% ==============================================================================
% Title and Authors
% ==============================================================================
\title{\textbf{\Huge XPower Banq}\\[0.5em]
\Large A Permissionless DeFi Lending Protocol with\\Lethargic Governance and Beta-Distributed Position Caps}

\author{
    \banqAuthor\thanks{\banqAuthorContact}
}

\date{}

% ==============================================================================
% Document Begin
% ==============================================================================
\begin{document}

\ifbundle\else
\maketitle
% ==============================================================================
% arxiv-tag.tex
% The arXiv-style vertical preprint tag printed on the first page of each
% standalone paper and on the bundle title page. The date is taken from
% \banqArxivDate (defined in banq-meta.tex) and may be overridden locally
% before % ==============================================================================
% arxiv-tag.tex
% The arXiv-style vertical preprint tag printed on the first page of each
% standalone paper and on the bundle title page. The date is taken from
% \banqArxivDate (defined in banq-meta.tex) and may be overridden locally
% before % ==============================================================================
% arxiv-tag.tex
% The arXiv-style vertical preprint tag printed on the first page of each
% standalone paper and on the bundle title page. The date is taken from
% \banqArxivDate (defined in banq-meta.tex) and may be overridden locally
% before \input{../shared/arxiv-tag} via \renewcommand.
% ==============================================================================

\begin{tikzpicture}[remember picture, overlay]
  \node[rotate=90, anchor=center, font=\Huge, color=codegray]
    at ([xshift=1cm]current page.west)
    {\banqArxivId~~\banqArxivDate};
\end{tikzpicture}
 via \renewcommand.
% ==============================================================================

\begin{tikzpicture}[remember picture, overlay]
  \node[rotate=90, anchor=center, font=\Huge, color=codegray]
    at ([xshift=1cm]current page.west)
    {\banqArxivId~~\banqArxivDate};
\end{tikzpicture}
 via \renewcommand.
% ==============================================================================

\begin{tikzpicture}[remember picture, overlay]
  \node[rotate=90, anchor=center, font=\Huge, color=codegray]
    at ([xshift=1cm]current page.west)
    {\banqArxivId~~\banqArxivDate};
\end{tikzpicture}

\fi

% ==============================================================================
% Abstract
% ==============================================================================
\banqAbstract{%
We present \textbf{XPower Banq}, a permissionless DeFi lending protocol on the Ethereum Virtual Machine. The protocol introduces: (1)~\emph{optionally locked positions} that attenuate liquidation cascades by factor~$(1{-}\phi)$, (2)~\emph{lethargic governance} with time-weighted parameter transitions bounded by $0.5\times$--$2\times$ per cycle, (3)~\emph{beta-distributed position caps} with holder-count scaling that rate-limits capacity accumulation to $O(\sqrt{k})$ for $k$~accounts, (4)~\emph{transferable debt positions} with inverted ERC20 semantics enabling capital-efficient debt assumption liquidation, and (5)~\emph{log-space TWAP oracles} with bidirectional geometric mean spread computation and logarithmic spread scaling. The protocol balances capital efficiency against bounded bad-debt risk: default parameters yield 66.67\% LTV with a 50\% over-collateralization buffer, adjustable via lethargic governance (\S\ref{subsec:lethargic-governance}). We provide technical analysis of the architecture, security properties, simulation results, and limitations.%
}{DeFi, Lending Protocol, Liquidation, Oracle, TWAP, Governance, Interest Rates}

% ==============================================================================
% 1. Introduction
% ==============================================================================
\banqPaper{Protocol Whitepaper}{A Permissionless DeFi Lending Protocol with Lethargic Governance and Beta-Distributed Position Caps}

\section{Introduction}
\label{sec:introduction}

Decentralized lending protocols---Compound~\cite{compound2019}, Aave~\cite{aave2020}, MakerDAO~\cite{makerdao2017}---have established algorithmic interest rate markets as foundational DeFi infrastructure~\cite{bartoletti2021}. Persistent challenges remain: oracle manipulation~\cite{mackinga2022}, governance attacks~\cite{beanstalk2022}, capital inefficiency, and systemic liquidation cascades~\cite{gudgeon2020}.

XPower Banq addresses these through five design choices:

\begin{enumerate}[leftmargin=*]
\item \textbf{Locked Positions}: Irrevocable position locks attenuate liquidation cascades by factor~$(1{-}\phi)$ through preventing immediate collateral redemption (Section~\ref{subsec:locked-positions}).

\item \textbf{Lethargic Governance}: Multiplicative bounds ($0.5\times$--$2\times$) with asymptotic transitions prevent governance shocks (Section~\ref{subsec:lethargic-governance}).

\item \textbf{Beta-Distributed Caps}: The \mbox{$12\lambda(1{-}\lambda)^2/\sqrt{n{+}2}$} cap function rate-limits capacity accumulation sublinearly (Section~\ref{subsec:beta-caps}).

\item \textbf{Debt Assumption Liquidation}: Transferable debt positions with inverted ERC20 semantics eliminate liquidator capital requirements (Section~\ref{subsec:health-liquidation}).

\item \textbf{Log-Space TWAP Oracle}: Geometric mean temporal averaging with bidirectional spread computation resists flash loan manipulation (Section~\ref{subsec:oracle-twap}).
\end{enumerate}

The remainder of this paper is organized as follows. Section~\ref{sec:related} surveys related work and positions our contributions against existing protocols. Section~\ref{sec:architecture} presents the contract architecture, followed by Section~\ref{sec:mechanisms} which details the core mechanisms. Section~\ref{sec:antispam} describes anti-spam protection, and Section~\ref{sec:governance} covers the governance model. Section~\ref{sec:security} provides a security analysis under an explicit adversary model. Section~\ref{sec:evaluation} presents simulation and gas evaluation results, Section~\ref{sec:limitations} discusses limitations, and Section~\ref{sec:conclusion} concludes. Appendices~A--C provide specifications and reference tables. A companion document~\cite{mtp,sim,ref} provides extended mathematical analysis, formal proofs, simulation implementations, and reference material.

% ==============================================================================
% 2. Related Work
% ==============================================================================
\section{Related Work}
\label{sec:related}

\textbf{Lending Protocols.} Compound~\cite{compound2019} introduced algorithmic interest rate markets with utilization-based rate curves and 2-day governance timelocks. Aave~\cite{aave2020} extended this with flash loans, stable rate borrowing, and credit delegation; V3 introduced E-Mode for correlated assets, though liquidation still requires liquid capital. MakerDAO~\cite{makerdao2017} pioneered the CDP model, suffering \$8.3M in losses during ``Black Thursday''~\cite{gudgeon2020}. Euler~\cite{euler2022} pioneered permissionless asset listing with risk-tiered collateral and a reactive rate model. Liquity~\cite{liquity2021} introduced the Stability Pool where pre-deposited capital absorbs liquidations; XPower Banq instead transfers debt to liquidators. Morpho~\cite{morpho2022} optimizes rates via peer-to-peer matching atop existing protocols.

\textbf{Liquidation and MEV.} Gudgeon et al.~\cite{gudgeon2020} analyzed the March~2020 crisis in which cascade effects drained significant protocol value. Perez et al.~\cite{perez2021} showed that liquidation markets are highly concentrated, with the top~5 liquidators capturing over 80\% of value, and Daian et al.~\cite{daian2020} and Qin et al.~\cite{qin2022} quantified MEV extraction from liquidation events.

\textbf{Oracles.} Mackinga et al.~\cite{mackinga2022} demonstrated practical TWAP manipulation, motivating several mitigation approaches. Chainlink~\cite{chainlink2017} provides off-chain aggregation but carries staleness risks during network congestion~\cite{werner2022}. Uniswap~V3~\cite{uniswap2021} introduced geometric mean TWAP computed in ring buffers, and Angeris and Chitra~\cite{angeris2020} formalized manipulation cost bounds for constant function market makers.

\textbf{Governance.} Beanstalk's \$182M governance attack~\cite{beanstalk2022} exploited same-block execution, underscoring the need for rate-bounded governance. Standard timelocks~\cite{openzeppelin2020} constrain the timing of changes but not their magnitude. XPower Banq addresses this with lethargic governance that requires multiple cycles for any significant parameter change.

\textbf{Comparative Analysis.} Table~\ref{tab:protocol-comparison} compares protocols across key dimensions. Some features marked absent may be deliberately omitted by design (e.g., Morpho Blue delegates oracle validation to curators), and XPower Banq's advantages carry trade-offs discussed in Section~\ref{sec:limitations}.

\begin{table*}[htbp]
\centering
\caption{Comparative analysis of DeFi lending protocols. Entries reflect architectural choices; absence of a feature may be deliberate.}
\label{tab:protocol-comparison}
\small
\begin{tabular}{@{}lcccccc@{}}
\toprule
Feature & Compound & Aave & MakerDAO & Liquity & Euler & \textbf{XPower Banq} \\
\midrule
1. Liquidation Model & Repayment & Repayment & Auction & Stability Pool & Repayment & Debt Assumption \\
2. Liquidator Capital & Required & Required & Required & Pre-deposited & Required & Not Required \\
3. Cascade Attenuation & None & None & Partial & Yes (pool) & None & Yes (locks)\textsuperscript{$\dagger$} \\
4. Position Transfer & Supply only & Supply only & No & No & Supply only & Both (inverted) \\
5. Governance Bounds & None & None & None & Immutable & None & $0.5\times$--$2\times$ \\
6. Param.\ Transitions & Instant & Instant & Instant & N/A & Instant & Asymptotic \\
7. Oracle Type & Chainlink & Chainlink & CL+OSM & Chainlink & Uniswap & Log TWAP \\
8. Capacity Rate-Limit & None & None & None & None & None & $\sqrt{n{+}2}$ scaling \\
9. Spam Protection & Gas price & Gas price & Gas price & Gas price & Gas price & Gas + PoW \\
10. Position Caps & None & Isolation & Debt ceil. & Debt ceil. & Borrow caps & Beta-distributed \\
\bottomrule
\end{tabular}
\vspace{0.3em}
\par\raggedright\scriptsize\textsuperscript{$\dagger$}Attenuation proportional to lock adoption $\phi$; without locks, cascade risk equals traditional protocols.
\end{table*}

% ==============================================================================
% 3. Protocol Architecture
% ==============================================================================
\section{Protocol Architecture}
\label{sec:architecture}

The protocol comprises six core contracts, depicted in Figure~\ref{fig:architecture}:

\begin{figure}[htbp]
\centering
\begin{tikzpicture}[
    node distance=1cm,
    box/.style={rectangle, draw, minimum width=1.8cm, minimum height=0.7cm, align=center, font=\scriptsize},
    arrow/.style={->, thick}
]
    % Core contracts
    \node[box, fill=primaryblue!20] (pool) {Pool};
    \node[box, fill=gray!20, above=0.8cm of pool] (acma) {Acma};
    \node[box, fill=orange!20, right=0.8cm of pool] (oracle) {Oracle};
    \node[box, fill=successgreen!20, below left=0.8cm and 0.3cm of pool] (supply) {Supply\\Position};
    \node[box, fill=warningred!20, below right=0.8cm and 0.3cm of pool] (borrow) {Borrow\\Position};
    \node[box, fill=yellow!20, below=1.8cm of pool] (vault) {Vault};

    % Arrows
    \draw[arrow] (acma) -- (pool);
    \draw[arrow] (pool) -- (oracle);
    \draw[arrow] (pool) -- (supply);
    \draw[arrow] (pool) -- (borrow);
    \draw[arrow] (pool) -- (vault);
    \draw[arrow] (vault) -- (supply);
    \draw[arrow] (vault) -- (borrow);
\end{tikzpicture}
\caption{XPower Banq contract architecture}
\label{fig:architecture}
\end{figure}

\begin{enumerate}[leftmargin=*]
    \item \textbf{Pool}: The main contract, managing supply, borrow, settle, and redeem operations with health checks and liquidation logic.

    \item \textbf{Position}: ERC20~\cite{erc20} tokens representing supply and borrow positions, each with distinct transfer semantics (Section~\ref{subsec:transfer}).

    \item \textbf{Vault}: An ERC4626-compliant~\cite{erc4626} custody contract that tracks deposited assets and utilization.

    \item \textbf{Oracle}: A log-space TWAP aggregator with bidirectional geomean spread computation.

    \item \textbf{Acma}: An access-control contract extending OpenZeppelin's \texttt{AccessManager}~\cite{openzeppelin2020} with the protocol-specific role catalogue and the \texttt{relate()} selector-to-role binder.

    \item \textbf{WPosition}: Optional ERC20 wrappers for supply/borrow positions, registered per-token via \texttt{Pool.enwrap()}; enable composition with external DeFi protocols that expect plain ERC20 transfer semantics (see Section~\ref{sec:limitations}).
\end{enumerate}

\begin{definition}[Minimum Token Requirements]
\label{def:token-constraints}
$|\mathcal{T}| \geq 2 \land \text{decimals}(T_i) \geq 6$, ensuring sufficient precision for interest calculations and cross-collateralization.
\end{definition}

% ==============================================================================
% 4. Core Mechanisms
% ==============================================================================
\section{Core Mechanisms}
\label{sec:mechanisms}

\subsection{Locked Positions and Cascade Attenuation}
\label{subsec:locked-positions}

\begin{definition}[Position Lock]
\label{def:position-lock}
For each user $u$ and position type $p \in \{\text{supply}, \text{borrow}\}$, the protocol maintains $\text{balance}(u,p)$, $\text{lock}(u,p)$, and $\text{liquid}(u,p) = \text{balance} - \text{lock}$. Operations accept a term parameter \texttt{dt\_term}; when nonzero, supplied assets cannot be redeemed and borrowed amounts cannot be settled until the slot epoch elapses. The sentinel \texttt{dt\_term = type(uint256).max} produces a permanent (irrevocable) lock; any other term writes into one of 16 quarterly ring slots that expire automatically. Transfers proportionally move locked amounts:
\begin{equation}
\label{eq:lock-transfer}
\text{transfer}_{\text{lock}}(u_1, u_2, v) = \left\lfloor \frac{\text{lock}(u_1) \cdot v}{\text{balance}(u_1)} \right\rfloor
\end{equation}
\end{definition}

\textbf{Lock Terms.} Locks are \emph{time-bound by default}: callers pass a term \texttt{dt\_term} that pins the principal to one of 16 quarterly ring slots (default slot length \texttt{MAX\_LOCK\_TERM} $=$ 3~months, total horizon $\approx 48$ months). Slots expire automatically and are reclaimed by \texttt{free()}. Passing \texttt{dt\_term = type(uint256).max} upgrades the lock to \emph{permanent} (irrevocable); only this path matches the cascade-attenuation analysis below. Cascade attenuation by factor $(1{-}\phi)$ therefore applies only to the \emph{permanent-locked} fraction; timed locks contribute attenuation only until their slot epoch elapses. A user can otherwise exit only via transfer (which proportionally moves the lock) or liquidation. Crucially, only the \emph{principal} is locked: accrued interest on locked supply positions can be redeemed at any time, and accrued interest on locked borrow positions can be settled at any time, too. Still, the design creates a fundamental tension: cascade protection depends on widespread adoption, yet rational adoption requires holding periods exceeding 5~years at typical utilization \cite{mtp}. The protocol addresses this tension by tying APY incentives to utilization levels.

\textbf{Lock Bonus/Malus.} To encourage adoption, locked suppliers earn additional interest $I_{\text{bonus}} = I \cdot \rho_u \cdot r_{\text{bonus}}$, while locked borrowers pay reduced interest $I_{\text{malus}} = I \cdot \rho_u \cdot r_{\text{malus}}$, where $\rho_u = \text{lock}/\text{balance}$ and $r_{\text{bonus}}, r_{\text{malus}} \in [0, s]$. The code enforces $r_{\text{bonus}} \leq s$ and $r_{\text{malus}} \leq s$ independently, where $s$ is the spread parameter. At full lock adoption ($\bar{\rho} = 1$), the combined constraint $r_{\text{bonus}} + r_{\text{malus}} \leq 2s$ holds at the boundary, with protocol margin approaching zero.

\begin{theorem}[Cascade Attenuation]
\label{thm:cascade}
In a pool where fraction $\phi$ of supply positions are locked, the maximum cascade amplification factor is bounded by $(1{-}\phi)$. Locked positions attenuate sell pressure by preventing immediate redemption of seized collateral; they do not eliminate cascades among unlocked positions.
\end{theorem}

Because locked supply tokens cannot be redeemed, liquidation transfers position tokens without underlying assets leaving the vault. Only the unlocked fraction $(1{-}\phi)$ generates sell pressure. A full proof appears in \cite{mtp}. Simulation confirms that at a 25\% price shock, unlocked positions suffer 85.7\% liquidation versus 29.4\% when fully locked \cite{sim}.

\textbf{Liquidation Precedence.} In practice, liquidators prefer unlocked positions for their immediate liquidity, creating \emph{de facto} seniority for locked positions. This ordering is market-driven rather than protocol-enforced.

\subsection{Position Transfer Semantics}
\label{subsec:transfer}

Supply positions follow standard ERC20 semantics: \texttt{transfer} pushes value to the receiver, and the health check is applied to the sender. Borrow positions, by contrast, implement \emph{inverted} semantics that reflect the nature of debt: \texttt{transfer(from, amount)} \textbf{pulls} debt from the \texttt{from} address, \texttt{transferFrom} requires \emph{dual} approval from both sender and receiver, and the health check falls on the receiver.

\begin{figure}[H]
\centering
\begin{tikzpicture}[
    node distance=2.2cm,
    box/.style={rectangle, draw, rounded corners, minimum width=1.4cm, minimum height=0.9cm, align=center, font=\footnotesize},
    arrow/.style={->, >=stealth, thick},
    label/.style={font=\scriptsize, align=center}
]
\node[box] (caller1) {Caller\\(from)};
\node[box, right=of caller1] (receiver1) {Receiver\\(to)};
\draw[>-, >=stealth, thick, successgreen!70] (caller1) -- node[above, label] {\textbf{push supply}} (receiver1);
\node[below=0.05cm of caller1, font=\scriptsize, warningred!70] {\textbf{health check}};

\node[box, below=1.4cm of caller1] (from2) {From};
\node[box, right=of from2] (to2) {To};
\node[box, above=0.6cm of $(from2)!0.5!(to2)$] (caller2) {Caller};
\draw[arrow, codegray, dashed] (from2) -- node[left, label, xshift=-2pt] {approve} (caller2);
\draw[>-, >=stealth, thick, successgreen!70] (from2) -- node[above, label] {\textbf{push supply}} (to2);
\node[below=0.05cm of from2, font=\scriptsize, warningred!70] {\textbf{health check}};
\end{tikzpicture}
\caption{Supply position --- standard ERC20 push semantics}
\label{fig:supply-transfer}
\end{figure}

\begin{figure}[H]
\centering
\begin{tikzpicture}[
    node distance=2.2cm,
    box/.style={rectangle, draw, rounded corners, minimum width=1.4cm, minimum height=0.9cm, align=center, font=\footnotesize},
    arrow/.style={->, >=stealth, thick},
    label/.style={font=\scriptsize, align=center}
]
\node[box] (source1) {Source\\(from)};
\node[box, right=of source1] (caller1) {Caller\\(to)};
\draw[arrow, warningred!70] (source1) -- node[above, label] {\textbf{pull borrow}} (caller1);
\node[below=0.05cm of caller1, font=\scriptsize, warningred!70] {\textbf{health check}};

\node[box, below=1.4cm of source1] (from2) {From};
\node[box, right=of from2] (to2) {To};
\node[box, above=0.6cm of $(from2)!0.5!(to2)$] (caller2) {Caller};
\draw[arrow, codegray, dashed] (from2) -- node[left, label, xshift=-2pt] {approve} (caller2);
\draw[arrow, codegray, dashed] (to2) -- node[right, label, xshift=2pt] {approve} (caller2);
\draw[arrow, warningred!70] (from2) -- node[above, label] {\textbf{pull borrow}} (to2);
\node[below=0.05cm of to2, font=\scriptsize, warningred!70] {\textbf{health check}};
\end{tikzpicture}
\caption{Borrow position --- inverted pull semantics with dual approval}
\label{fig:borrow-transfer}
\end{figure}

\begin{definition}[Debt Transfer]
\label{def:debt-transfer}
When debt is transferred from user $A$ to user $B$: (1)~$A$'s borrow position decreases by $d$; (2)~$B$'s borrow position increases by $d$; (3)~$B$ must have sufficient collateral; (4)~both must approve the intermediary (unless self-transfer).
\end{definition}

\textbf{Rationale.} Debt transfer differs fundamentally from asset transfer: with assets, the receiver benefits; with debt, the sender benefits by losing an obligation. The inverted semantics reflect this distinction---standard ERC20 pushes value to the receiver, whereas borrow positions pull debt from the source. The Pool contract as owner bypasses approvals (and health checks) to enable liquidations.

\subsection{Lethargic Governance}
\label{subsec:lethargic-governance}

\begin{definition}[Parameter Transition]
\label{def:param-transition}
A parameter $\theta$ transitions from $\theta_0$ to $\theta_1$ subject to $0.5 \cdot \theta_0 \leq \theta_1 \leq 2 \cdot \theta_0$. The effective value follows the time-weighted mean:
\begin{equation}
\bar{\theta}(t) = \theta_1 - \frac{(\theta_1 - \theta_0)(t_s - t_0)}{t - t_0}
\end{equation}
which asymptotically approaches $\theta_1$ as $t \to \infty$.
\end{definition}

\begin{figure}[htbp]
\centering
\begin{tikzpicture}
\begin{axis}[
    width=\columnwidth,
    height=5cm,
    xlabel={Time},
    ylabel={Parameter Value},
    xmin=0, xmax=10,
    ymin=0.95, ymax=1.25,
    xtick={0,2,10},
    xticklabels={$t_0$,$t_s$,$\infty$},
    ytick={1.0,1.1,1.2},
    yticklabels={$\theta_0$,,\ $\theta_1$},
    grid=major,
    legend pos=south east,
]
\addplot[thick, primaryblue!70, domain=0:2, forget plot] {1.0};
\addplot[thick, primaryblue!70, domain=2:10, samples=50] {1.2 - 0.4/x};
\addplot[thick, dashed, warningred!70, domain=0:10] {1.0};
\addplot[thick, dashed, successgreen!70, domain=2:10] {1.2};
\addplot[only marks, mark=*, mark size=2pt, primaryblue!70, forget plot] coordinates {(2,1.0)};
\legend{Actual, Source, Target}
\end{axis}
\end{tikzpicture}
\caption{Lethargic parameter transition (asymptotic)}
\label{fig:lethargic-transition}
\end{figure}

\textbf{Safety Constraints.} Each single change is bounded to at most $2\times$, overlapping transitions are disallowed, and changes are rate-limited to once per governance period (monthly). Mandatory initial lock periods ranging from 3~months to 1~year (by category) prevent manipulation at deployment. Full parameter tables are provided in Appendix~\ref{apx:parameters}.

\subsection{Beta-Distributed Position Caps}
\label{subsec:beta-caps}

\begin{figure}[htbp]
\centering
\begin{tikzpicture}
\begin{axis}[
    width=\columnwidth,
    height=5cm,
    xlabel={$\lambda = B/S$ (Position Fraction)},
    ylabel={Beta Multiplier},
    xmin=0, xmax=1,
    ymin=0, ymax=1.8,
    xtick={0,0.33,0.5,1.0},
    xticklabels={0,1/3,1/2,1},
    grid=major,
    legend pos=north east,
]
\addplot[thick, primaryblue!70, domain=0.01:0.99, samples=100] {12*x*(1-x)^2};
\addplot[thick, dashed, warningred!70] coordinates {(0.333,0) (0.333,1.6)};
\addplot[only marks, mark=*, mark size=2pt, primaryblue!70, forget plot] coordinates {(0.333,1.778)};
\node at (axis cs:0.75,1.65) [anchor=center, font=\small] {Peak at $\lambda=1/3$};
\end{axis}
\end{tikzpicture}
\caption{Beta distribution cap function \mbox{$12\lambda(1{-}\lambda)^2$}}
\label{fig:beta-distribution}
\end{figure}

\begin{definition}[Position Cap Function]
\label{def:cap}
For balance $B$, supply $S$, and $n$ large holders ($\geq 1$ token unit):
\begin{equation}
\label{eq:beta-cap}
\text{cap}(B, S, n) = \frac{C_{\max} \cdot 12\lambda(1{-}\lambda)^2}{\sqrt{n + 2}}
\end{equation}
where $\lambda = B/S$ and $C_{\max}$ is the maximum \emph{relative} cap.
\end{definition}

The offset~$+2$ prevents degeneracy at low holder counts. In addition, governance sets a holder floor $n_{\min}$ via the \texttt{MIN\_HOLDERS} parameter; the effective count is $\max(n, n_{\min})$, which tightens per-user caps during the cold-start phase when few holders have joined. When $\lambda = 0$ (a new user with no existing balance), the beta factor is bypassed and the cap reduces to $C_{\max}/\sqrt{n{+}2}$, ensuring a non-zero entry allocation.

The Beta(2,3) shape is a design choice, not a uniquely optimal distribution---any unimodal density vanishing at 0 and 1 achieves similar goals. We chose Beta(2,3) for its mode at $\lambda = 1/3$ (favoring medium positions) and analytic tractability. A simpler $4\lambda(1{-}\lambda)$ (Beta(2,2)) or $\lambda(1{-}\lambda)$ would also penalize extremes; the asymmetry toward smaller positions is a preference, not a theoretical necessity.

\textbf{Rate-Limiting, Not Sybil Prevention.} The $\sqrt{n{+}2}$ divisor bounds the \emph{accumulation rate}: $k$ accounts yield $O(\sqrt{k})$ total capacity gain per iteration. This does not, however, penalize long-term equilibrium share---simulation shows that a Sybil attacker retains 48.6\% share (versus 50\% initial) after 60~weeks, a marginal 1.4\% penalty. The mechanism thus prevents rapid capacity monopolization while accepting that patient attackers with sufficient capital can reach similar long-term positions. Combined with a 7~iterations/week cap, capacity growth remains bounded regardless of account splitting. See \cite{sim} for detailed simulations.

\subsection{Health Factor and Liquidation}
\label{subsec:health-liquidation}

\begin{definition}[Health Factor]
\label{def:health}
For user $u$ with supply values $V_s^i$ and borrow values $V_b^i$ across $n$ tokens:
\begin{equation}
\label{eq:health}
H(u) = \frac{\sum_i w_s^i \cdot V_s^i}{\sum_i w_b^i \cdot V_b^i}
\end{equation}
where $w_s^i, w_b^i \in [0, 255]$ are uint8 weights. The implementation divides each weighted value by $\texttt{WEIGHT\_MAX} = 255$ and averages across tokens; both normalizations cancel in the ratio.
\end{definition}

With the default weights $w_s = 170$ and $w_b = 255$, the effective LTV is $170/255 \approx 66.67\%$. This is standard over-collateralization---the same mechanism underlying Compound's and Aave's LTV parameters, expressed through weight ratios rather than an explicit LTV setting. The resulting implicit liquidation bonus $\beta = w_b/w_s - 1 = 50\%$ incentivizes liquidators.

\textbf{Hybrid Liquidation.} The protocol supports both centralized keepers (the default, via \texttt{square()}) and governance-enabled public liquidators (PoW-gated, via \texttt{liquidate()}). Both paths use a \emph{debt assumption} model in which the liquidator assumes fraction $2^{-e}$ of the victim's positions. For a given exponent $e$, both supply tokens $s = \text{supply}(V) \gg e$ and debt $d = \text{borrow}(V) \gg e$ transfer atomically. No liquid capital is required---only sufficient collateral headroom. The liquidator's health is verified post-liquidation.

\subsection{Oracle TWAP}
\label{subsec:oracle-twap}

The oracle aggregates prices via a log-space exponential moving average (EMA). Each refresh queries the AMM in both directions (source$\to$target and target$\to$source) and computes a bidirectional geomean spread, which eliminates directional bias in asymmetric pools:
\begin{align}
\overline{m}_t &= \alpha \cdot \overline{m}_{t-1} + (1{-}\alpha) \cdot \log_2(\bar{p}_{t-1}) \\
\overline{r}_t &= \alpha \cdot \overline{r}_{t-1} + (1{-}\alpha) \cdot r_{t-1}
\end{align}
where $r_t = \frac{1}{2}(\log_2(1+s^{AB}_t) + \log_2(1+s^{BA}_t))$. Smoothing in $\log_2$ space produces geometric mean temporal averaging (as in Uniswap~V3~\cite{uniswap2021}). The implementation defers the most recent tick by one period (\texttt{twap\_.last $\to$ twap\_.mean} on the next refresh) to provide two-tick flash-loan immunity; see \texttt{library/TWAP.sol::update}. With $\alpha = 0.944$ (12-period half-life) and hourly refreshes, approximately 40~hours of sustained manipulation achieves 90\% price deviation.

\textbf{Spread Scaling.} Large positions incur wider spreads via $\mu = \log_2(2x + 2)$ where $x = n \cdot s$, reflecting market impact without requiring per-pair parameters. The base spread itself serves as sensitivity control: liquid pairs scale slowly, while illiquid pairs scale quickly.

\textbf{Bad-Debt Risk from Stale Pricing.} The slow convergence is double-edged. After a genuine $2\times$ price move, only 22.5\% is absorbed after 14~hours. Combined with the 1-hour refresh limit, the protocol can be 2+ hours blind during a crash. If collateral drops 50\%, the oracle reports near pre-crash prices for hours, potentially delaying liquidations and allowing bad debt to accumulate. The 50\% over-collateralization buffer provides margin: at the default 66.67\% LTV the protocol sits on the boundary where bad debt begins to emerge under tail-event price crashes ($>$50\%); Monte Carlo simulation and analytical bounds quantify the resulting risk profile. The conservative governance floor of 33\% LTV retains zero bounded bad debt for crashes up to 50\%. See \cite{sim} for TWAP simulations and quantitative bad-debt risk bounds.

\subsection{Interest Rates}
\label{subsec:interest-rates}

\begin{figure}[htbp]
\centering
\begin{tikzpicture}
\begin{axis}[
    width=\columnwidth,
    height=5.5cm,
    xlabel={Utilization $U$},
    ylabel={Interest Rate $R$},
    xmin=0, xmax=1,
    ymin=0, ymax=1.1,
    xtick={0,0.9,1.0},
    xticklabels={0,$U^*$,1},
    ytick={0,0.1,0.5,1.0},
    yticklabels={0,$R^*$,0.5,1},
    grid=major,
    legend pos=north west,
]
\addplot[thick, primaryblue!70, domain=0:0.9, forget plot] {x/9};
\addplot[thick, primaryblue!70, domain=0.9:1.0, forget plot] {9*x - 8};
\addplot[thick, dashed, gray!50, forget plot] coordinates {(0.9,0) (0.9,0.1)};
\addplot[only marks, mark=*, primaryblue!70] coordinates {(0.9,0.1)};
\addlegendentry{Kink point}
\node at (axis cs:0.45,0.05) [anchor=south, font=\small, rotate=4] {Linear};
\node at (axis cs:0.92,0.55) [anchor=center, font=\small, rotate=83] {Steep};
\end{axis}
\end{tikzpicture}
\caption{Piecewise-linear interest rate curve}
\label{fig:interest-curve}
\end{figure}

\begin{definition}[Interest Rate Model]
\label{def:irm}
For utilization $U$, optimal utilization $U^*$, and optimal rate $R^*$:
\begin{equation}
\label{eq:interest-rate}
\begin{aligned}
R(U) &= \min\!\left(R_{\text{kink}}(U),\ 2\right), \\
R_{\text{kink}}(U) &= \begin{cases}
U R^* / U^* & \text{if } U \leq U^* \\[4pt]
\dfrac{U(1{-}R^*) - (U^* - R^*)}{1{-}U^*} & \text{otherwise}
\end{cases}
\end{aligned}
\end{equation}
The implementation caps the base rate at 200\% to bound interest in the post-kink regime.
\end{definition}

The protocol maintains a symmetric spread: $R_{\text{borrow}} = R_{\text{base}} \cdot (1{+}s)$, $R_{\text{supply}} = R_{\text{base}} \cdot (1{-}s)$, generating $2s$ margin on interest flows. See \cite{mtp} for mathematical foundations including interest accrual and time-weighted integration.

% ==============================================================================
% 5. Anti-Spam Protection
% ==============================================================================
\section{Anti-Spam Protection}
\label{sec:antispam}

When governance enables public operations---such as liquidation via \texttt{liquidate()} or oracle refresh via \texttt{refresh()}---proof-of-work gates prevent spam. The PoW requires:
\begin{equation}
\text{zeros}\bigl(\text{keccak256}(\text{blockHash} \| \text{tx.origin} \| \text{msg.data})\bigr) \geq d
\end{equation}
where each difficulty level requires $16\times$ more computational work. Combined with token-bucket rate limiting (capacity $C$, regeneration 1/second), this prevents mempool flooding and block stuffing. By default, operations are restricted to authorized keepers; PoW only applies in governance-enabled public mode.

\textbf{Fairness Limitations.} PoW shifts the advantage from capital (gas wars) to computation. Professional miners achieve 10--20$\times$ higher hash rates than browser-based implementations, and on PoS chains (e.g., Avalanche) validators who know they will propose a block can pre-compute favorable nonces, creating a validator-advantage vector. In short, PoW provides anti-spam friction rather than fair access---it changes the extraction mechanism from capital to computation without eliminating it.

% ==============================================================================
% 6. Governance and Parameters
% ==============================================================================
\section{Governance and Parameters}
\label{sec:governance}

All protocol parameters are governed via the lethargic transition mechanism (Definition~\ref{def:param-transition}). They are organized into four contract domains: Pool (rate limits, PoW difficulty, health weights), Position (interest rate model, caps), Oracle (TWAP behavior, feed management), and Vault (entry/exit fees). Full parameter tables and role categories appear in Appendix~\ref{apx:parameters}.

\begin{theorem}[Governance Rate Bound]
\label{thm:governance}
Under multiplicative bound $\mu = 2$ and minimum cycle $\Delta t_{\min}$, the parameter at time $t$ satisfies:
\begin{equation}
\mu^{-n} \theta_0 \leq \theta(t) \leq \mu^n \theta_0
\end{equation}
where $n = \lfloor(t{-}t_0)/\Delta t_{\min}\rfloor$. A $10\times$ change requires $\lceil\log_2 10\rceil = 4$ months.
\end{theorem}

The protocol implements a three-tier role hierarchy per governable function: an executor (who calls the function), an admin (who grants and revokes the executor role), and a guard (who provides veto-only revocation). Guard roles serve as a ``break glass'' mechanism, allowing emergency halts of malicious actions without granting execution power. The governance rate bound proof is provided in \cite{mtp}.

% ==============================================================================
% 7. Security Analysis
% ==============================================================================
\section{Security Analysis}
\label{sec:security}

\subsection{Adversary Model}

We consider an adversary $\mathcal{A}$ who controls capital $K$, hash rate $H$, arbitrarily many accounts, governance access with probability $p_g$, and same-block front-running capability. We assume that $\mathcal{A}$ cannot break cryptographic primitives or control more than 50\% of consensus.

\subsection{Core Security Properties}

\textbf{Cascade Attenuation (Theorem~\ref{thm:cascade}).} Because locked supply positions cannot be redeemed, cascade sell pressure is bounded to $(1{-}\phi)$ of the unlocked pool. This attenuates---but does not prevent---cascades. Simulation confirms that at a 25\% price shock, unlocked positions suffer 85.7\% liquidation versus 29.4\% when fully locked (Section~\ref{sec:evaluation}).

\textbf{Rate-Limited Accumulation.} The $\sqrt{n{+}2}$ divisor bounds per-iteration capacity gain sublinearly. Combined with 7~iterations/week, reaching 99\% capacity requires approximately $850\sqrt{n/100}$ iterations. This constitutes rate-limiting rather than Sybil prevention: the long-term share distribution is ultimately determined by initial capital \cite{sim}.

\textbf{Governance Rate Bound (Theorem~\ref{thm:governance}).} Catastrophic parameter changes require 6+ months of sustained malicious governance, providing ample time for detection and response.

\textbf{Solvency.} The code enforces $r_{\text{bonus}} \leq s$ and $r_{\text{malus}} \leq s$ independently per parameter. At full lock adoption, the combined constraint $r_{\text{bonus}} + r_{\text{malus}} \leq 2s$ holds at the boundary, with margin approaching zero. The aggregate solvency formula ($r_{\text{bonus}} \cdot \bar{\rho}^S + r_{\text{malus}} \cdot \bar{\rho}^B \leq 2s$) is a sufficient condition that is not directly enforced on-chain; the per-parameter bounds serve as the operative constraints.

\subsection{Risk Assessment}

\begin{table}[htbp]
\centering
\caption{Risk severity assessment}
\label{tab:risks}
\small
\begin{tabularx}{0.95\columnwidth}{@{}l >{\centering\arraybackslash}X >{\centering\arraybackslash}X@{}}
\toprule
Risk Category & Likelihood & Impact \\
\midrule
Oracle Manipulation & Medium & High \\
Governance Attack & Low & High \\
Liquidation Cascade\textsuperscript{$\dagger$} & Medium & High \\
Smart Contract Bug & Medium & Critical \\
Keeper Centralization & Medium & High \\
Bad Debt (Extreme Vol.) & Low & Critical \\
Lock Mechanism Abuse & Low & Medium \\
PoW Resource Advantage & Medium & Low \\
Oracle Staleness & Medium & Medium \\
\bottomrule
\end{tabularx}
\vspace{0.5em}
\par\raggedright\scriptsize\textsuperscript{$\dagger$}Impact conditional on lock adoption; without locks, impact is High.
\end{table}

\textbf{Oracle.} The log-space TWAP with $\alpha = 0.944$ requires approximately 40~hours of sustained manipulation to achieve 90\% deviation, and flash loans are ineffective due to the two-tick immunity window. However, the 2+~hour blindness during genuine crashes creates bad-debt risk (Section~\ref{subsec:oracle-twap}).

\textbf{Governance.} Parameter changes are bounded by $0.5\times$--$2\times$ per cycle, and guard roles enable emergency revocation. Nevertheless, an attacker who sustains control for 6+~months can achieve significant cumulative drift.

\textbf{Smart Contracts.} The implementation uses reentrancy protection via \texttt{ReentrancyGuardTransient}, Solidity 0.8+ overflow checks, and OpenZeppelin access control. No formal verification has been performed (Section~\ref{sec:limitations}).

Full proofs of the above properties are provided in \cite{mtp}.

% ==============================================================================
% 8. Evaluation
% ==============================================================================
\section{Evaluation}
\label{sec:evaluation}

\subsection{Agent-Based Cascade Simulation}

We simulate 1,000 borrowers with log-normal positions, health factors $H \sim \mathcal{N}(1.5, 0.3)$ truncated at $[1.0, 3.0]$, and linear price impact $\Delta p = k \cdot V_{\text{sold}}$ with $k = 1.5 \times 10^{-4}$ (pool ${\sim}25\%$ of market depth).

\begin{table}[htbp]
\centering
\caption{Cascade simulation --- positions liquidated (\%), $k = 1.5 \times 10^{-4}$}
\label{tab:cascade-sim}
\small
\begin{tabularx}{0.95\columnwidth}{@{}>{\raggedleft\arraybackslash}X >{\raggedleft\arraybackslash}X >{\raggedleft\arraybackslash}X >{\raggedleft\arraybackslash}X >{\raggedleft\arraybackslash}X >{\raggedleft\arraybackslash}X@{}}
\toprule
Lock $\phi$ & 10\% & 15\% & 20\% & 25\% & 30\% \\
\midrule
  0\% & 8.6 & 19.8 & 36.5 & 85.7 & 99.1 \\
 25\% & 8.2 & 16.5 & 30.6 & 55.8 & 84.9 \\
 50\% & 7.8 & 14.4 & 26.7 & 37.1 & 61.2 \\
 75\% & 7.8 & 13.5 & 22.4 & 33.3 & 47.9 \\
100\% & 7.4 & 12.7 & 19.8 & 29.4 & 38.6 \\
\bottomrule
\end{tabularx}
\end{table}

\begin{figure}[htbp]
\centering
\begin{tikzpicture}
\begin{axis}[
    width=\columnwidth,
    height=5.5cm,
    xlabel={Price Drop (\%)},
    ylabel={Positions Liquidated (\%)},
    xmin=10, xmax=30,
    ymin=0, ymax=100,
    xtick={10,15,20,25,30},
    grid=major,
    legend pos=north west,
    legend style={font=\scriptsize},
    legend cell align=right,
]
\addplot[thick, warningred!70, mark=*] coordinates {(10,8.6) (15,19.8) (20,36.5) (25,85.7) (30,99.1)};
\addplot[thick, orange!80, mark=square*] coordinates {(10,8.2) (15,16.5) (20,30.6) (25,55.8) (30,84.9)};
\addplot[thick, yellow!90!orange, mark=triangle*] coordinates {(10,7.8) (15,14.4) (20,26.7) (25,37.1) (30,61.2)};
\addplot[thick, successgreen!70, mark=diamond*] coordinates {(10,7.8) (15,13.5) (20,22.4) (25,33.3) (30,47.9)};
\addplot[thick, primaryblue!70, mark=pentagon*] coordinates {(10,7.4) (15,12.7) (20,19.8) (25,29.4) (30,38.6)};
\legend{$\phi=0\%$, $\phi=25\%$, $\phi=50\%$, $\phi=75\%$, $\phi=100\%$}
\end{axis}
\end{tikzpicture}
\caption{Cascade depth vs.\ price drop for varying lock fractions}
\label{fig:cascade-sim}
\end{figure}

\textbf{Simulation Limitations.} The simulation uses 1,000 agents, which is too few for high statistical confidence. Health factor distributions are assumed rather than calibrated against real DeFi data. The linear price impact model~\cite{kyle1985continuous} is a first-order approximation; real order books exhibit concave impact. No confidence intervals or sensitivity analysis on $k$ are provided, and historical stress events are discussed narratively but not backtested. Implementation details appear in \cite{sim}.

\subsection{Gas Cost Analysis}

Gas measurements are post-optimization Foundry snapshots (\texttt{FOUNDRY\_OPTIMIZER=true}, default profile, \texttt{via\_ir=false}), benchmarked against Aave V3.3 figures from a 2025 Cyfrin gas-optimization audit. Banq's warm-path gas is competitive with mainstream lending protocols on supply and dominates them on every other operation:

\begin{table}[htbp]
\centering
\caption{Banq vs.\ Aave V3 warm-path gas (Foundry, optimizer ON)}
\label{tab:gas-comparison}
\small
\begin{tabularx}{0.95\columnwidth}{@{}X rrr@{}}
\toprule
Operation & Banq & Aave V3 & $\Delta$ \\
\midrule
Supply             & 167{,}209 & 146{,}354 & $+$14\% \\
Borrow             & 140{,}884 & 247{,}485 & $-$43\% \\
Repay (settle)     & 123{,}403 & 189{,}518 & $-$35\% \\
Withdraw (redeem)  & 128{,}380 & 181{,}430 & $-$29\% \\
Liquidation (full) & 298{,}882 & 389{,}059 & $-$23\% \\
\bottomrule
\end{tabularx}
\end{table}

\begin{table}[htbp]
\centering
\caption{Cold-vs-warm and Banq-unique operations}
\label{tab:gas-cold-warm}
\small
\begin{tabularx}{0.95\columnwidth}{@{}X rrl@{}}
\toprule
Operation & Cold & Warm & Notes \\
\midrule
Supply             & 285{,}814 & 167{,}209 & --- \\
Borrow             & 244{,}833 & 140{,}884 & self-pair \\
Borrow (cross-pair) & ---       & 194{,}317 & +53k oracle leg \\
Liquidation (full) & 298{,}882 & ---       & --- \\
healthOf (view)    & 44{,}611  & ---       & view-only \\
lockSupply (timed)  & 179{,}860 & 66{,}688 & ring-slot lock \\
lockSupply (perma)  & 86{,}712  & ---      & irrevocable \\
lockBorrow (timed)  & 180{,}904 & 67{,}720 & debt lock \\
xfer\_supply (1 slot) & 191{,}463 & --- & lock-aware ERC20 \\
xfer\_supply (16 slots) & 402{,}525 & --- & worst-case ring scan \\
\bottomrule
\end{tabularx}
\end{table}

The savings come from two implementation choices: (i)~the log-space compounding index (\cite{log}) eliminates \texttt{exp()} from the write path, saving $\sim$1{,}200~gas per accrual at the cost of $\sim$1{,}100~gas per \texttt{totalOf} read; (ii)~the packed \texttt{\_state} (global) and \texttt{\_stateOf} (per-user) storage words pack \texttt{[uint80~index | uint64~stamp/depth | uint112~balance]} into a single \texttt{SSTORE}, halving the write cost relative to the canonical three-slot layout. The \texttt{Lock} library applies the same technique: \texttt{uint128[16]} ring slots pack two epoch-value pairs per word (8~words instead of 16~per user), and a single \texttt{cache} word packs \texttt{[uint120~perma | uint120~total | uint16~bits]}. On Avalanche (approximately 1.55~gwei, AVAX~$<$\$10), all operations cost well under 1~cent~USD.

% ==============================================================================
% 9. Limitations and Future Work
% ==============================================================================
\section{Limitations and Future Work}
\label{sec:limitations}

\textbf{Capital Efficiency.} The default 66.67\% LTV requires \$1.50 of collateral per \$1 borrowed---comparable to Compound's 75\% and within 25~percentage points of Aave V3 E-Mode (93\%). Lethargic governance (\S\ref{subsec:lethargic-governance}) bounds each cycle to $0.5\times$--$2\times$ the prior target, so the default is reachable down to a 33\% conservative floor in one cycle, and the bad-debt analysis in \cite{sim} characterises the risk profile across the LTV configurations swept therein; no analysis is provided of what LTV range the target market (XPower tokens on Avalanche) actually requires.

\textbf{Formal Verification.} The smart contracts have not undergone formal verification. Given the 10~interacting mechanisms---where locks affect caps, which affect health, which affects liquidation---this represents a significant gap that undermines the security analysis. Formal proofs of health factor correctness, position token conservation, and interest accrual monotonicity are absent. This is the highest-priority future work, recommended via Certora, Halmos, or the K~framework.

\textbf{Composability Risks.} Wrapped positions (WPosition) interact with external DeFi protocols, yet no analysis addresses how inverted borrow semantics or irrevocable locks behave under composition. External contracts that expect standard ERC20 semantics may mishandle borrow positions.

\textbf{Missing Game-Theoretic Analysis.} No model addresses competitive liquidator dynamics under PoW constraints, equilibrium liquidation profit, or oracle sandwich attack quantification. While PoW raises the cost of sandwich attacks, it does not eliminate them.

\textbf{Future Directions.} Areas for further development include dynamic parameter adjustment based on market conditions, cross-chain lending via message passing, privacy-preserving positions via zero-knowledge proofs, and integration of additional oracle sources for off-chain price anchoring.

% ==============================================================================
% 10. Conclusion
% ==============================================================================
\section{Conclusion}
\label{sec:conclusion}

XPower Banq presents a DeFi lending protocol that balances capital efficiency against bounded bad-debt risk. Its strongest contributions are lethargic governance---a genuine innovation with formal rate bounds---and the log-space geometric-mean TWAP oracle, which is both novel and empirically validated. The locked position mechanism provides cascade attenuation proportional to adoption, though its effectiveness depends on rational adoption dynamics that remain uncertain \cite{mtp}. The beta-distributed cap function rate-limits capacity accumulation but does not prevent long-term Sybil advantage. The default 66.67\% LTV positions the protocol within range of mainstream lending markets while remaining adjustable via lethargic governance (\S\ref{subsec:lethargic-governance}). Formal verification and game-theoretic liquidation modeling remain the most significant gaps for future work.

% ==============================================================================
% References
% ==============================================================================
% ==============================================================================
% Appendices (standalone only; in the bundle these continue as plain
% sections of Part~I to keep counters consistent across all subfiles).
% ==============================================================================
\ifbundle\else
\appendix
\fi

\section{Protocol Parameters}
\label{apx:parameters}

\subsection{Pool Parameters}

\begin{table}[H]
\centering
\caption{Pool parameters. Time units: s=second, d=day, w=week, m=month, y=year.}
\label{tab:pool-params}
\small
\begin{tabularx}{0.95\columnwidth}{@{}l l@{\,}c@{\,}r X@{}}
\toprule
Parameter & \multicolumn{3}{c}{Setting} & Description \\
\midrule
\texttt{MAX\_SUPPLY} & 1w & $\in$ & [1s, 1y] & Max supply rate-limit \\
\texttt{MIN\_SUPPLY} & 1d & $\in$ & [1s, 1y] & Min supply rate-limit \\
\texttt{POW\_SUPPLY} & 0 & $\in$ & [0, 64] & Supply PoW difficulty \\
\texttt{MAX\_BORROW} & 1w & $\in$ & [1s, 1y] & Max borrow rate-limit \\
\texttt{MIN\_BORROW} & 1d & $\in$ & [1s, 1y] & Min borrow rate-limit \\
\texttt{POW\_BORROW} & 0 & $\in$ & [0, 64] & Borrow PoW difficulty \\
\texttt{POW\_SQUARE} & 0 & $\in$ & [0, 64] & Liquid. PoW difficulty \\
\texttt{WEIGHT\_SUPPLY} & 170 & $\in$ & [0, 255] & Supply health weight\textsuperscript{$\ddagger$} \\
\texttt{WEIGHT\_BORROW} & 255 & $\in$ & [0, 255] & Borrow health weight \\
\bottomrule
\end{tabularx}
\vspace{0.3em}
\par\raggedright\scriptsize\textsuperscript{$\ddagger$}Deployer-set; reference deployments use 170/255 $\approx$ 66.67\% LTV.
\end{table}

\subsection{Position Parameters}

\begin{table}[H]
\centering
\caption{Position parameters}
\label{tab:position-params}
\small
\begin{tabularx}{0.95\columnwidth}{@{}l l@{\,}c@{\,}r X@{}}
\toprule
Parameter & \multicolumn{3}{c}{Setting} & Description \\
\midrule
\texttt{UTIL} & 90\% & $\in$ & [0, 100\%] & Optimal utilization \\
\texttt{RATE} & 10\% & $\in$ & [0, 100\%] & Rate at kink \\
\texttt{SPREAD} & 10\% & $\in$ & [0, 50\%] & Rate half-spread\textsuperscript{$\S$} \\
\texttt{LOCK\_BONUS} & 10\% & $\in$ & [0, \texttt{SPREAD}] & Locked supply APY bonus \\
\texttt{LOCK\_MALUS} & 10\% & $\in$ & [0, \texttt{SPREAD}] & Locked borrow APY reduction \\
\texttt{CAP} & --- & $\in$ & [0, $2^{112}{-}1$] & Position cap ($\approx 5.19 \times 10^{15}$ at 18 decimals) \\
\texttt{MIN\_HOLDERS} & --- & $\in$ & [0, $10^{18}$] & Min large holders \\
\bottomrule
\end{tabularx}
\vspace{0.3em}
\par\raggedright\scriptsize\textsuperscript{$\S$}Lower-bounded by \texttt{LOCK\_BONUS} and \texttt{LOCK\_MALUS}.
\end{table}

\subsection{Oracle Parameters}

\begin{table}[H]
\centering
\caption{Oracle parameters}
\label{tab:oracle-params}
\small
\begin{tabularx}{0.95\columnwidth}{@{}l l@{\,}c@{\,}r X@{}}
\toprule
Parameter & \multicolumn{3}{c}{Setting} & Description \\
\midrule
\texttt{DECAY} & 0.944 & $\in$ & [0.5, 1.0] & EMA decay \\
\texttt{LIMIT} & 1h & $\in$ & [1s, 1d] & Min refresh interval \\
\texttt{LEVEL} & 0 & $\in$ & [0, 64] & Refresh PoW difficulty \\
\texttt{DELAY} & 14d & $\in$ & [1w, 3m] & Feed enlist timelock \\
\bottomrule
\end{tabularx}
\end{table}

\subsection{Vault Parameters}

\begin{table}[H]
\centering
\caption{Vault parameters}
\label{tab:vault-params}
\small
\begin{tabularx}{0.95\columnwidth}{@{}l l@{\,}c@{\,}r X@{}}
\toprule
Parameter & \multicolumn{3}{c}{Setting} & Description \\
\midrule
\texttt{FEE\_ENTRY} & 0.1\% & $\in$ & [0, 50\%] & Deposit fee \\
\texttt{FEE\_EXIT} & 1.0\% & $\in$ & [0, 50\%] & Withdrawal fee \\
\bottomrule
\end{tabularx}
\end{table}

\subsection{Change Rate Constraints}

\begin{table}[H]
\centering
\caption{Initial lock periods of parameters}
\label{tab:transition-durations}
\small
\begin{tabularx}{0.95\columnwidth}{@{}X >{\raggedleft\arraybackslash}X@{}}
\toprule
Category & Lock Period \\
\midrule
Interest rate model & 3m \\
Oracle TWAP settings & 3m \\
Oracle feed delay & 1y \\
Vault fees & 3m \\
Pool weights & 3m \\
Pool rate limits & 1y \\
\bottomrule
\end{tabularx}
\end{table}

\begin{table}[H]
\centering
\caption{Maximum parameter change over time}
\label{tab:max-change}
\small
\begin{tabularx}{0.95\columnwidth}{@{}>{\centering\arraybackslash}X r r@{}}
\toprule
Cycles & Max Increase & Max Decrease \\
\midrule
1 & $2\times$ & $0.5\times$ \\
2 & $4\times$ & $0.25\times$ \\
3 & $8\times$ & $0.125\times$ \\
6 & $64\times$ & $0.016\times$ \\
12 & $4096\times$ & $0.00024\times$ \\
\bottomrule
\end{tabularx}
\end{table}

\subsection{Role-Based Access Control}

The protocol implements a three-tier role hierarchy per governable function: executor (calls the restricted function), admin (grants/revokes the executor role), and guard (veto-only revocation).

\textbf{Governance Attack Precautions.} Delayed feed changes require up to 3~months. Guard role separation prevents escalation. Gradual parameter migration eliminates arbitrage from sudden changes. Bounded cumulative change creates predictable worst-case scenarios. Immutable constraints include minimum decimals $\geq 6$, minimum 2~tokens, a fixed token list, and non-upgradeable contracts.

\section{Protocol Constants}
\label{apx:constants}

The following tables list all compile-time constants defined in the protocol's \texttt{Constant.sol} library.

\begin{table}[H]
\centering
\caption{Time unit constants (in seconds)}
\label{tab:time-constants}
\small
\begin{tabularx}{0.95\columnwidth}{@{}l r X@{}}
\toprule
Constant & Value & Definition \\
\midrule
\texttt{CENTURY} & 3,155,760,000 & $365.25 \times 100 \times 86400$ \\
\texttt{YEAR} & 31,557,600 & $365.25 \times 86400$ \\
\texttt{MONTH} & 2,629,800 & YEAR / 12 \\
\texttt{WEEK} & 604,800 & $7 \times 86400$ \\
\texttt{DAY} & 86,400 & $24 \times 3600$ \\
\texttt{HOUR} & 3,600 & $60 \times 60$ \\
\texttt{MINUTE} & 60 & 60 seconds \\
\texttt{SECOND} & 1 & 1 second \\
\bottomrule
\end{tabularx}
\end{table}

\begin{table}[H]
\centering
\caption{Fixed-point representation constants}
\label{tab:fixedpoint-constants}
\small
\begin{tabularx}{0.95\columnwidth}{@{}l r X@{}}
\toprule
Constant & Value & Meaning \\
\midrule
\texttt{ONE} & $10^{18}$ & WAD unit (1.0) \\
\texttt{TWO} & $2 \times 10^{18}$ & WAD unit (2.0) \\
\texttt{HLF} & $0.5 \times 10^{18}$ & WAD unit (0.5) \\
\texttt{NIL} & 0 & WAD unit (0.0) \\
\texttt{PCT} & $10^{16}$ & 1 percentage point \\
\texttt{BPS} & $10^{14}$ & 1 basis point \\
\bottomrule
\end{tabularx}
\end{table}

\begin{table}[H]
\centering
\caption{Protocol constraints and limits}
\label{tab:constraint-constants}
\small
\begin{tabularx}{0.95\columnwidth}{@{}l r X@{}}
\toprule
Constant & Value & Purpose \\
\midrule
\texttt{MAX\_DIFFICULTY} & 64 & Max PoW difficulty \\
\texttt{MIN\_HOLDERS} & $10^{18}$ & Min holders upper bound \\
\texttt{VERSION} & 0x1ba & Protocol version \\
\bottomrule
\end{tabularx}
\end{table}

\section{EMA Decay Factors}
\label{apx:decay}

Table~\ref{tab:decay} lists pre-computed EMA decay factors relating half-life (in refresh periods) to the decay coefficient $\alpha$, used for log-space TWAP oracle smoothing.

\begin{table}[H]
\centering
\caption{Pre-computed EMA decay factors}
\label{tab:decay}
\small
\begin{tabularx}{0.95\columnwidth}{@{}>{\centering\arraybackslash}X >{\raggedleft\arraybackslash}X >{\centering\arraybackslash}X@{}}
\toprule
Half-life & Decay $\alpha$ & Use Case \\
\midrule
1 & 0.500000 & Fast response \\
2 & 0.707107 & Short-term \\
12 & 0.943874 & Medium-term \\
24 & 0.971532 & Long-term \\
\bottomrule
\end{tabularx}
\end{table}

\ifbundle\else
\begin{thebibliography}{99}

% === Cited references, in first-citation order ===

% §1 Introduction (line 132)
\bibitem{compound2019}
Leshner, R. and Hayes, G.
\newblock Compound: The Money Market Protocol.
\newblock \emph{Compound Whitepaper}, 2019.
\newblock \url{https://compound.finance/documents/Compound.Whitepaper.pdf}

\bibitem{aave2020}
Aave Protocol.
\newblock Aave Protocol Whitepaper V2.0.
\newblock \emph{Aave Documentation}, 2020.
\newblock \url{https://docs.aave.com/}

\bibitem{makerdao2017}
MakerDAO Team.
\newblock The Maker Protocol: MakerDAO's Multi-Collateral Dai (MCD) System.
\newblock \emph{MakerDAO Whitepaper}, 2017.
\newblock \url{https://makerdao.com/whitepaper/}

\bibitem{bartoletti2021}
Bartoletti, M., Chiang, J.H., and Lafuente, A.L.
\newblock SoK: Lending Pools in Decentralized Finance.
\newblock \emph{Financial Cryptography and Data Security (FC)}, LNCS vol.~12676, pp.~553--578, 2021.

\bibitem{mackinga2022}
Mackinga, T., Nadahalli, T., and Wattenhofer, R.
\newblock TWAP Oracle Attacks: Easier Done than Said?
\newblock \emph{IEEE International Conference on Blockchain and Cryptocurrency}, 2022.

\bibitem{beanstalk2022}
Immunefi.
\newblock Hack Analysis: Beanstalk Governance Attack, April 2022.
\newblock \emph{Security Report}, 2022.
\newblock \url{https://medium.com/immunefi/hack-analysis-beanstalk-governance-attack-april-2022-f42788fc821e}

\bibitem{gudgeon2020}
Gudgeon, L., Perez, D., Harz, D., Livshits, B., and Gervais, A.
\newblock The Decentralized Financial Crisis.
\newblock \emph{Crypto Valley Conference}, 2020.

% §2 Related Work — Lending (line 156)
\bibitem{euler2022}
Euler Labs.
\newblock Euler Finance: Permissionless Lending Protocol.
\newblock \emph{Euler Whitepaper}, 2022.
\newblock \url{https://docs.euler.finance/}

\bibitem{liquity2021}
Lauko, R. and Pardoe, R.
\newblock Liquity: Decentralized Borrowing Protocol.
\newblock \emph{Liquity Whitepaper}, 2021.
\newblock \url{https://docs.liquity.org/}

\bibitem{morpho2022}
Frambot, P. and Gontier Delaunay, M.
\newblock Morpho Optimizer: Optimizing Decentralized Liquidity Protocols.
\newblock \emph{Morpho Whitepaper}, 2022.
\newblock \url{https://whitepaper.morpho.org/}

% §2 Related Work — Liquidation and MEV (line 158)
\bibitem{perez2021}
Perez, D., Werner, S.M., Xu, J., and Livshits, B.
\newblock Liquidations: DeFi on a Knife-edge.
\newblock \emph{Financial Cryptography and Data Security}, 2021.

\bibitem{daian2020}
Daian, P., Goldfeder, S., Kell, T., Li, Y., Zhao, X., Bentov, I., Breidenbach, L., and Juels, A.
\newblock Flash Boys 2.0: Frontrunning in Decentralized Exchanges, Miner Extractable Value, and Consensus Instability.
\newblock \emph{IEEE Symposium on Security and Privacy (S\&P)}, pp.~910--927, 2020.

\bibitem{qin2022}
Qin, K., Zhou, L., and Gervais, A.
\newblock Quantifying Blockchain Extractable Value: How dark is the forest?
\newblock \emph{IEEE Symposium on Security and Privacy (S\&P)}, pp.~198--214, 2022.

% §2 Related Work — Oracles (line 160)
\bibitem{chainlink2017}
Ellis, S., Juels, A., and Nazarov, S.
\newblock ChainLink: A Decentralized Oracle Network.
\newblock \emph{Chainlink Whitepaper}, 2017.
\newblock \url{https://chain.link/whitepaper}

\bibitem{werner2022}
Werner, S.M., Perez, D., Gudgeon, L., Klages-Mundt, A., Harz, D., and Knottenbelt, W.J.
\newblock SoK: Decentralized Finance (DeFi).
\newblock \emph{ACM Computing Surveys}, 2022.

\bibitem{uniswap2021}
Adams, H., Zinsmeister, N., Salem, M., Keefer, R., and Robinson, D.
\newblock Uniswap v3 Core.
\newblock \emph{Uniswap Documentation}, 2021.
\newblock \url{https://docs.uniswap.org/}

\bibitem{angeris2020}
Angeris, G. and Chitra, T.
\newblock Improved Price Oracles: Constant Function Market Makers.
\newblock \emph{ACM Conference on Advances in Financial Technologies (AFT)}, 2020.

% §2 Related Work — Governance (line 162)
\bibitem{openzeppelin2020}
OpenZeppelin Team.
\newblock OpenZeppelin Contracts.
\newblock \emph{OpenZeppelin Documentation}, 2020.
\newblock \url{https://docs.openzeppelin.com/contracts/}

% §3 Architecture (lines 205, 207)
\bibitem{erc20}
Vogelsteller, F. and Buterin, V.
\newblock EIP-20: Token Standard.
\newblock \emph{Ethereum Improvement Proposals}, 2015.
\newblock \url{https://eips.ethereum.org/EIPS/eip-20}

\bibitem{erc4626}
Santoro, J., t11s, Jadeja, J., and Cuesta Ca\~{n}ada, A.
\newblock EIP-4626: Tokenized Vaults.
\newblock \emph{Ethereum Improvement Proposals}, 2022.
\newblock \url{https://eips.ethereum.org/EIPS/eip-4626}

% §8 Evaluation (line 561)
\bibitem{kyle1985continuous}
Kyle, A.S.
\newblock Continuous Auctions and Insider Trading.
\newblock \emph{Econometrica}, 53(6):1315--1335, 1985.

% === Uncited references (retained for further reading) ===

\bibitem{uniswap2020}
Adams, H., Zinsmeister, N., and Robinson, D.
\newblock Uniswap v2 Core.
\newblock \emph{Uniswap Documentation}, 2020.
\newblock \url{https://docs.uniswap.org/}

\bibitem{glassnode2020}
Glassnode Insights.
\newblock What Really Happened to MakerDAO on Black Thursday.
\newblock \emph{Glassnode Analysis}, March 2020.
\newblock \url{https://insights.glassnode.com/what-really-happened-to-makerdao/}

\bibitem{liu2023terra}
Liu, J., Makarov, I., and Schoar, A.
\newblock Anatomy of a Run: The Terra Luna Crash.
\newblock \emph{Harvard Law School Forum on Corporate Governance}, May 2023.
\newblock \url{https://corpgov.law.harvard.edu/2023/05/22/anatomy-of-a-run-the-terra-luna-crash/}

\bibitem{qin2021flashloans}
Qin, K., Zhou, L., Livshits, B., and Gervais, A.
\newblock Attacking the DeFi Ecosystem with Flash Loans for Fun and Profit.
\newblock \emph{Financial Cryptography and Data Security (FC)}, LNCS vol.~12674, pp.~3--32, 2021.

\bibitem{nakamoto2008}
Nakamoto, S.
\newblock Bitcoin: A Peer-to-Peer Electronic Cash System.
\newblock 2008.
\newblock \url{https://bitcoin.org/bitcoin.pdf}

\bibitem{douceur2002}
Douceur, J.R.
\newblock The Sybil Attack.
\newblock \emph{International Workshop on Peer-to-Peer Systems (IPTPS)}, LNCS vol.~2429, pp.~251--260, 2002.

\bibitem{feist2019}
Feist, J., Grieco, G., and Groce, A.
\newblock Slither: A Static Analysis Framework for Smart Contracts.
\newblock \emph{IEEE/ACM International Workshop on Emerging Trends in Software Engineering for Blockchain (WETSEB)}, pp.~8--15, 2019.

\bibitem{dwork1992}
Dwork, C. and Naor, M.
\newblock Pricing via Processing or Combatting Junk Mail.
\newblock \emph{Advances in Cryptology -- CRYPTO '92}, LNCS vol.~740, pp.~139--147, 1992.

\bibitem{roughgarden2021}
Roughgarden, T.
\newblock Transaction Fee Mechanism Design.
\newblock \emph{ACM Conference on Economics and Computation (EC)}, 2021.

\bibitem{log}
\textsc{Karun The Ritch}.
\newblock XPower Banq: Log-Space Compounding Index.
\newblock \emph{Preprint}, 2026.
\newblock \url{https://www.xpowerbanq.com}

\bibitem{mtp}
\textsc{Karun The Ritch}.
\newblock XPower Banq: Mathematical Theory \& Proofs.
\newblock \emph{Preprint}, 2026.
\newblock \url{https://www.xpowerbanq.com}

\bibitem{sim}
\textsc{Karun The Ritch}.
\newblock XPower Banq: Simulations \& Risk Analysis.
\newblock \emph{Preprint}, 2026.
\newblock \url{https://www.xpowerbanq.com}

\bibitem{ref}
\textsc{Karun The Ritch}.
\newblock XPower Banq: References \& Glossary.
\newblock \emph{Preprint}, 2026.
\newblock \url{https://www.xpowerbanq.com}

\end{thebibliography}
\fi

% ==============================================================================
% Glossary
% ==============================================================================
\ifbundle\else
\section*{Glossary}
\label{sec:glossary}
\addcontentsline{toc}{section}{Glossary}

\noindent A comprehensive glossary with over 120 entries appears in \cite{ref}. The following terms are central to this paper.

\begin{description}[leftmargin=!,labelwidth=2.5cm]
    \item[Beta Cap] Position limit following \mbox{$12\lambda(1-\lambda)^2$} distribution where $\lambda$ is the user's fraction of total supply; peaks at $\lambda = 1/3$ and vanishes at boundaries.

    \item[Borrow Position] ERC20 token representing debt obligation; uses inverted transfer semantics where \texttt{transfer} pulls debt.

    \item[Debt Assumption] Liquidation model where the liquidator assumes the victim's debt rather than repaying it; enables capital-efficient liquidation without requiring liquid capital.

    \item[Health Factor] Ratio of weighted supply value to weighted borrow value: $H = \sum w_s V_s / \sum w_b V_b$. Liquidation occurs when $H < 1$.

    \item[Inverted Transfer] Borrow position transfer semantics where \texttt{transfer(from, amount)} pulls debt FROM the first parameter rather than pushing to it. See \S4.2.

    \item[Kink] Utilization threshold (e.g., 90\%) where interest rate slope increases; incentivizes liquidity retention.

    \item[Lethargic Governance] Governance model with time-delayed, bounded parameter transitions. Values approach targets asymptotically, bounded to $0.5\times$--$2\times$ per governance cycle. See \S A.

    \item[Liquidation] Forced closure of an unhealthy position ($H < 1$); XPower Banq uses debt assumption model.

    \item[LTV] Loan-to-Value ratio---maximum borrowing power as fraction of collateral value; default $w_s/w_b = 170/255 \approx 66.67\%$.

    \item[Position Lock] Fraction $\phi$ of a position restricted from redemption or sale; prevents liquidation cascades.

    \item[Spread] (1)~Oracle: bidirectional geomean of rel.\ spreads from both AMM query directions, stored as $\log_2(1 + s_{\text{geo}})$; wider spreads indicate lower liquidity or higher manipulation resistance. (2)~IRM: symmetric half-spread parameter (e.g., 10\%) applied to base rate; borrow rate = base~$\times$~$(1{+}s)$, supply rate = base~$\times$~$(1{-}s)$.

    \item[Supply Position] ERC20 token representing deposited collateral; uses standard transfer semantics.

    \item[TWAP] Time-Weighted Average Price---price smoothed over time via log-space EMA (geometric mean temporal averaging) to resist manipulation. Simulated in \cite{sim}.

    \item[Utilization] Ratio of borrowed assets to supplied assets in a vault; drives interest rates via a piecewise-linear model with a kink at optimal utilization.

    \item[Vault] ERC4626-compliant custody contract holding deposited assets; tracks utilization for the interest rate model.
\end{description}

\clearpage
\ifodd\value{page}\else\thispagestyle{empty}\null\clearpage\fi
\fi

\end{document}
